\chapter{Model}
\label{chap:methodology}

\section{Motivation and overview}

The exploratory analyses presented in Chapter~\ref{chap:background} establish that the player--season feature space contains meaningful lower-dimensional structure. Principal component analysis revealed that approximately nineteen components are needed to account for ninety per cent of the total sample variance, while an eleven-factor oblimin solution provided an interpretable account of the empirical correlation structure through recognisable football dimensions such as goal threat, chance creation, pass circulation, defensive work, and duel efficiency. Classical multidimensional scaling applied to the cosine-based dissimilarity confirmed that the pairwise similarity structure admits a low-dimensional Euclidean approximation. Taken together, these findings support the view that player--season performance profiles are governed by a moderate number of shared latent traits, and that the feature space is rich enough for such traits to be recoverable from the data.

What the exploratory analyses do not do, however, is commit to a structure for how a player's statistics arise. PCA scores and factor loadings summarise the sample as given, treating every player--season row as exchangeable and saying nothing about why two rows from the same player should resemble each other more than two rows from different players. The model developed in this chapter states that structure directly: an additive decomposition that separates a stable player effect from a season effect, with a residual specification chosen to match how the data actually behave rather than assumed by default. This matters especially in a longitudinal setting, where the same player is observed in multiple seasons and two such observations are not independent: they share the player's stable latent characteristics, which the exploratory analyses' exchangeability assumption does not recognise.

This chapter specifies a probabilistic model designed to address these limitations. The model is hierarchical: each player's observed performance vector in a given season is the sum of three parts; a stable player effect, a season-specific shift, and residual noise. The player effects are themselves modelled as draws from a distribution whose covariance is structured as a low-rank factor model, directly encoding the finding from Chapter~\ref{chap:background} that player-level variation is largely explained by a small number of shared latent dimensions.

The model has three principal objectives. First, it provides a principled variance decomposition, attributing each feature's total variance to player, season, and residual components; this decomposition is expressed as a set of intraclass correlation coefficients, one per feature. Second, it estimates the covariance structure of the player effects through a factor model with $K$ shared latent dimensions, recovering interpretable directions in the player-effect space. Third, all inferences are Bayesian: the model produces a joint posterior distribution over all parameters, from which uncertainty in the latent structure, the variance decomposition, and player-level factor scores is directly quantifiable \cite{carpenter2017stan}.

\section{Notation}

Throughout this chapter, the following notation is used consistently. Let $N$ denote the total number of player--season observations, $I$ the number of unique players, $S$ the number of seasons, $P$ the number of standardised performance features, and $K$ a positive integer referred to as the factor rank, whose selection will be discussed in Section~\ref{sec:rank}. Observation indices run from $n = 1, \ldots, N$; $i[n] \in \{1, \ldots, I\}$ and $j[n] \in \{1, \ldots, S\}$ are the player and season indices of observation $n$, respectively.

\begin{table}[tbp]
\centering
\caption{Summary of notation used throughout this chapter.}
\label{tab:notation_model}
\begin{tabularx}{\textwidth}{lX}
\toprule
Symbol & Meaning \\
\midrule
$N,\, I,\, S,\, P,\, K$ & Total observations; unique players; seasons; features; factor rank \\
$n = 1,\ldots,N$ & Observation index \\
$i[n] \in \{1,\ldots,I\}$ & Player index of observation $n$ \\
$j[n] \in \{1,\ldots,S\}$ & Season index of observation $n$ \\
$\mathbf{y}_n \in \mathbb{R}^P$ & Standardised performance vector for observation $n$ \\
$\mathbf{a}_i \in \mathbb{R}^P$ & Player effect for player $i$ \\
$\mathbf{b}_j \in \mathbb{R}^P$ & Season effect for season $j$ \\
$\boldsymbol{\varepsilon}_n \in \mathbb{R}^P$ & Residual for observation $n$ \\
$\boldsymbol{\Lambda} \in \mathbb{R}^{P \times K}$ & Factor loading matrix \\
$\boldsymbol{\Psi}_a = \mathrm{diag}(\psi_{a,1}^2, \ldots, \psi_{a,P}^2)$ & Uniqueness matrix \\
$\boldsymbol{\Sigma}_a,\, \boldsymbol{\Sigma}_b,\, \boldsymbol{\Sigma}_e$ & Player, season, and residual covariance matrices \\
$\boldsymbol{\eta}_i \in \mathbb{R}^K$ & Latent factor scores for player $i$ \\
$\mathbf{z}_{u,i} \in \mathbb{R}^P$ & Latent uniqueness scores for player $i$ \\
$\nu > 2$ & Degrees-of-freedom parameter of the Student-$t$ residual \\
$\phi \geq 0$ & Observation-level variance scaling exponent \\
$m_n$ & Minutes played in observation $n$ \\
$m_{\mathrm{ref}}$ & Reference exposure level \\
\bottomrule
\end{tabularx}
\end{table}

\section{The additive hierarchical decomposition}

The central structural assumption of the model is that the performance vector for any given player--season observation arises from an additive decomposition into three independent latent components.

\begin{definition}[Additive decomposition]
\label{def:additive}
For each observation $n = 1, \ldots, N$, the observed performance vector $\mathbf{y}_n \in \mathbb{R}^P$ satisfies
\begin{equation}
\mathbf{y}_n = \mathbf{a}_{i[n]} + \mathbf{b}_{j[n]} + \boldsymbol{\varepsilon}_n,
\label{eq:additive}
\end{equation}
where $\mathbf{a}_i \in \mathbb{R}^P$ is the \emph{player effect} for player $i$, $\mathbf{b}_j \in \mathbb{R}^P$ is the \emph{season effect} for season $j$, and $\boldsymbol{\varepsilon}_n \in \mathbb{R}^P$ is the \emph{residual} for observation $n$. The three components are assumed mutually independent.
\end{definition}

The player effect $\mathbf{a}_i$ captures those aspects of player performance that are stable across seasons: the passing volume, defensive involvement, and goal-threat profile that characterise a player's identity rather than a particular team context. The season effect $\mathbf{b}_j$ captures shifts in the distribution of performance statistics that are common to all players in season $j$: changes in the pace of play, officiating style, or data-collection conventions that may change the observed values of many features simultaneously. The residual $\boldsymbol{\varepsilon}_n$ captures what remains after both systematic components are removed, including player--season idiosyncrasy, noise from the finite number of matches played, and any tactical or physical changes specific to one season that are not representative of the player's profile.

\begin{remark}[Absence of a global mean]
No explicit population mean vector $\boldsymbol{\mu}$ appears in \eqref{eq:additive}. This is because the features $\mathbf{y}_n$ are constructed by standardising each feature to zero mean and unit variance across the full sample, as described in Chapter~\ref{chap:data_source}. The population mean is therefore absorbed into the standardisation step, and the model is centred from the outset.
\end{remark}

The decomposition \eqref{eq:additive} defines a \emph{crossed} random-effects structure: players and seasons are crossed, not nested, each player appears in several seasons, and each season contains many players. This is in contrast to a nested structure in which seasons would be defined within players. The crossed structure is appropriate here because the inferential goal is to separate player-level persistence from season-level context, not to model players as a random sample within fixed seasons \cite{gelman2007data}.

\section{Covariance structure}
\label{sec:covariance}

The three components $\mathbf{a}_i$, $\mathbf{b}_j$, and $\boldsymbol{\varepsilon}_n$ are each assigned a zero-mean multivariate distribution; they differ only in the structure of their covariance matrices. The choices here for said components differ: a richer covariance structure is assigned to the player effects, which are the primary object of interest and for which the data contain many observations per latent quantity, while simpler structures are assigned to season effects and residuals.

\subsection{Player covariance: low-rank plus diagonal}

\begin{definition}[Factor covariance for player effects]
\label{def:sigma_a}
The player effects $\mathbf{a}_i$ are exchangeable with zero mean and covariance
\begin{equation}
\boldsymbol{\Sigma}_a = \boldsymbol{\Lambda}\boldsymbol{\Lambda}^\top + \boldsymbol{\Psi}_a,
\label{eq:sigma_a}
\end{equation}
where $\boldsymbol{\Lambda} \in \mathbb{R}^{P \times K}$ is the loading matrix and $\boldsymbol{\Psi}_a = \mathrm{diag}(\psi_{a,1}^2, \ldots, \psi_{a,P}^2)$ is the uniqueness matrix.
\end{definition}

The structure in \eqref{eq:sigma_a} decomposes the player-effect covariance into two parts. The term $\boldsymbol{\Lambda}\boldsymbol{\Lambda}^\top$ is a rank-$K$ positive semi-definite matrix that captures the shared dependence among features through $K$ latent dimensions; $\boldsymbol{\Psi}_a$ is a diagonal matrix that captures the feature-specific variance not explained by those shared dimensions.

\begin{proposition}[Entries of $\boldsymbol{\Sigma}_a$]
\label{prop:sigma_entries}
The diagonal entry $[\boldsymbol{\Sigma}_a]_{pp}$ for feature $p$ is
\begin{equation}
[\boldsymbol{\Sigma}_a]_{pp} = \sum_{k=1}^{K} \Lambda_{pk}^2 + \psi_{a,p}^2,
\label{eq:sigma_diag}
\end{equation}
and the off-diagonal entry for features $p \neq q$ is
\begin{equation}
[\boldsymbol{\Sigma}_a]_{pq} = \sum_{k=1}^{K} \Lambda_{pk}\,\Lambda_{qk}.
\label{eq:sigma_offdiag}
\end{equation}
\end{proposition}
\begin{proof}
Both statements follow by expanding $[\boldsymbol{\Lambda}\boldsymbol{\Lambda}^\top + \boldsymbol{\Psi}_a]_{pq}$. For the diagonal, $[\boldsymbol{\Lambda}\boldsymbol{\Lambda}^\top]_{pp} = \sum_k \Lambda_{pk}^2$ and $[\boldsymbol{\Psi}_a]_{pp} = \psi_{a,p}^2$. For the off-diagonal, $[\boldsymbol{\Lambda}\boldsymbol{\Lambda}^\top]_{pq} = \sum_k \Lambda_{pk}\Lambda_{qk}$ and $[\boldsymbol{\Psi}_a]_{pq} = 0$ since $\boldsymbol{\Psi}_a$ is diagonal.
\end{proof}

In the factor-analysis literature, the quantity $\sum_{k=1}^K \Lambda_{pk}^2$ is called the \emph{communality} of feature $p$: it measures how much of the player-level variance in feature $p$ is explained by the $K$ shared latent dimensions. The quantity $\psi_{a,p}^2$ is called the \emph{uniqueness}: it is the portion of player-level variance in feature $p$ that is specific to that feature and not captured by the shared structure. Together, they account for all player-level variance in feature $p$, as stated in \eqref{eq:sigma_diag}.

The off-diagonal entries in \eqref{eq:sigma_offdiag} reflect the pairwise covariance between features $p$ and $q$ induced by the shared factors: two features co-vary at the player level only to the extent that they share common loadings on the $K$ dimensions. A player who scores highly on a latent factor tends to score highly on all features with positive loadings on that factor, creating positive covariance among those features. The uniqueness matrix $\boldsymbol{\Psi}_a$ contributes nothing to off-diagonal entries, reflecting the assumption that, conditional on the $K$ shared factors, feature-specific residuals are uncorrelated.

This covariance structure is motivated directly by the exploratory findings of Chapter~\ref{chap:background}. There, classical factor analysis showed that the 48 features are not independent at the player level: they cluster into correlated groups, and a number of shared latent traits (eleven, by parallel analysis) were needed to reproduce their correlation pattern.

The low-rank form $\boldsymbol{\Sigma}_a = \boldsymbol{\Lambda}\boldsymbol{\Lambda}^\top + \boldsymbol{\Psi}_a$ builds that same idea into the model's assumptions. It says exactly this: player-level covariance among features is mostly explained by $K$ shared dimensions ($\boldsymbol{\Lambda}\boldsymbol{\Lambda}^\top$), plus a small, purely feature-specific leftover ($\boldsymbol{\Psi}_a$).

The difference from the earlier exploratory chapter is that this is now a fully probabilistic (Bayesian) model, not a one-off descriptive result \cite{lopes2004bayesian}. Classical factor analysis gives a single point-estimate set of loadings, with no sense of how uncertain they are. Here, $\boldsymbol{\Lambda}$, $\boldsymbol{\Psi}_a$, and everything built from them instead get a full posterior distribution: the model does not just claim that shared structure exists, it quantifies how confident it is about it. The rank $K$ adopted for $\boldsymbol{\Sigma}_a$ is selected separately in Section~\ref{sec:rank}.

\subsection{Season covariance: diagonal}

\begin{remark}[Diagonal season covariance]
\label{rem:sigma_b}
The season effects $\mathbf{b}_j$ are assigned zero mean and diagonal covariance:
\begin{equation}
\boldsymbol{\Sigma}_b = \mathrm{diag}(\sigma_{b,1}^2, \ldots, \sigma_{b,P}^2).
\label{eq:sigma_b}
\end{equation}
\end{remark}

This is an asymmetric design choice relative to the player covariance. A richer specification, $\boldsymbol{\Sigma}_b = \boldsymbol{\Lambda}_b\boldsymbol{\Lambda}_b^\top + \boldsymbol{\Psi}_b$, would model correlated season-level drift, but would require reliable estimation of $P \times K_b$ loading parameters from only $S$ observations of the season. With $S = 12$ seasons, this is severely under-determined for any $K_b \geq 1$: preliminary fits with a low-rank $\boldsymbol{\Sigma}_b$ produced R-hat values exceeding $1.70$ and bulk effective sample sizes as low as $6$ for the season loadings, indicating that the posterior was entirely dominated by the prior. The diagonal specification \eqref{eq:sigma_b} allows each feature its own season-level drift magnitude while avoiding parameters that the data cannot identify, plus the computational convenience.

\subsection{Residual covariance}

The residuals $\boldsymbol{\varepsilon}_n$ are assigned a diagonal covariance $\boldsymbol{\Sigma}_e = \mathrm{diag}(\sigma_{e,1}^2, \ldots, \sigma_{e,P}^2)$, reflecting the assumption that, conditional on the player and season effects, the $P$ features are uncorrelated at the observation level. The diagonal entries themselves are not constant across observations in the production model: Section~\ref{sec:phi} makes $\sigma_{e,p}$ depend on the observation's playing time, $\sigma_{e,n,p}$, to account for the fact that per-90 features are noisier estimates for players with fewer minutes. The cross-feature independence encoded by the diagonal structure is unaffected by that refinement.

\section{Likelihood specification}
\label{sec:likelihood}

The additive decomposition \eqref{eq:additive} and the covariance structure of Section~\ref{sec:covariance} are common to all model variants considered here. The models differ in the distributional form assigned to the scalar residuals $\varepsilon_{np}$, the $p$-th component of $\boldsymbol{\varepsilon}_n$.

\subsection{Normal baseline}

The simplest complete specification assigns each residual component a Gaussian distribution:
\begin{equation}
\varepsilon_{np} \mid \sigma_{e,p} \;\sim\; \mathcal{N}(0,\, \sigma_{e,p}^2), \qquad p = 1, \ldots, P,
\label{eq:normal_resid}
\end{equation}
independently across features and observations, giving the conditional joint likelihood
\begin{equation}
\mathbf{y}_n \mid \mathbf{a}_{i[n]},\, \mathbf{b}_{j[n]},\, \boldsymbol{\sigma}_e \;\sim\; \mathcal{N}_P\!\left(\mathbf{a}_{i[n]} + \mathbf{b}_{j[n]},\; \boldsymbol{\Sigma}_e\right).
\label{eq:normal_likelihood}
\end{equation}
This specification, referred to as the Normal model, serves as a computational and inferential baseline against which richer likelihood choices are compared.

\subsection{Student-$t$ residuals}
\label{sec:t_likelihood}

For player--season data, the Gaussian assumption on residuals is restrictive. Some player--seasons produce extreme per-90 statistics for identifiable reasons: a player returning from injury may appear in only a handful of matches with atypical physical load; a player deployed in an unusual tactical role for one season may display statistics far from any stable average; a breakout season may represent an extreme but transient realisation. Under the Normal model, these extreme residuals exert a disproportionate influence on the estimates of the player and season effects, since the Gaussian density penalises departures from the conditional mean quadratically and without limit. A heavier-tailed distribution assigns such observations a higher residual likelihood and effectively down-weights their influence on the shared parameters.

\begin{definition}[Student-$t$ error model]
\label{def:student_t}
In the Student-$t$ model, the residual for feature $p$ in observation $n$ follows
\begin{equation}
\varepsilon_{np} \mid \nu,\, \sigma_{e,p} \;\sim\; t\!\left(\nu,\; 0,\; \sigma_{e,p}\right), \qquad p = 1, \ldots, P,
\label{eq:t_resid}
\end{equation}
independently across features and observations, where $\nu > 2$ is a degrees-of-freedom parameter shared across all features and observations, and $\sigma_{e,p} > 0$ is the feature-specific scale. The Student-$t$ distribution with $\nu$ degrees of freedom, location zero, and scale $\sigma$ has density
\begin{equation}
f(\varepsilon;\, \nu,\, \sigma) \;=\; \frac{\Gamma\!\left(\tfrac{\nu+1}{2}\right)}{\Gamma\!\left(\tfrac{\nu}{2}\right)\sqrt{\nu\pi}\;\sigma} \left(1 + \frac{\varepsilon^2}{\nu\sigma^2}\right)^{\!-(\nu+1)/2}.
\label{eq:t_density}
\end{equation}
\end{definition}

\begin{proposition}[Variance of the Student-$t$ residual]
\label{prop:t_variance}
Under Definition~\ref{def:student_t}, the variance of $\varepsilon_{np}$ is
\begin{equation}
\mathrm{Var}(\varepsilon_{np}) \;=\; \sigma_{e,p}^2\;\frac{\nu}{\nu - 2},
\label{eq:t_variance}
\end{equation}
which is finite for $\nu > 2$ and converges to $\sigma_{e,p}^2$ as $\nu \to \infty$.
\end{proposition}
\begin{proof}
The variance of a Student-$t$ random variable with $\nu$ degrees of freedom and scale $\sigma$ is $\sigma^2 \nu/(\nu - 2)$, valid for $\nu > 2$. As $\nu \to \infty$, $\nu/(\nu-2) \to 1$.
\end{proof}

The tail-weight parameter $\nu$ has a direct interpretation: small values of $\nu$ produce heavier tails and greater likelihood of extreme residuals; large values recover the Gaussian behaviour \cite{lange1989robust,geweke1993bayesian}. The degrees-of-freedom parameter is treated as unknown and estimated from the data; its posterior summarises the evidence in favour of heavy-tailed residuals conditional on the model specification and the prior.

Proposition~\ref{prop:t_variance} assumes the constant residual scale $\sigma_{e,p}$ of this section's base specification. Section~\ref{sec:phi} replaces $\sigma_{e,p}$ by an observation-dependent scale in the production model, and the residual variance below is restated accordingly (Corollary~\ref{cor:t_variance_mv}).

\subsection{Observation-level variance scaling}
\label{sec:phi}

The residual scale $\sigma_{e,p}$ in \eqref{eq:normal_resid} and \eqref{eq:t_resid} is constant across observations. Section~\ref{sec:t_likelihood}'s Student-$t$ specification is stated this way because it is the natural point to introduce heavy tails on their own, but a constant residual scale does not hold for this data and is not what the production model uses. The performance features $\mathbf{y}_n$ are constructed as per-90 rates or efficiency ratios, both of which are ratio estimators: for a player who accumulated $m_n$ minutes, the sampling variance of the per-90 estimate of feature $p$ is approximately proportional to $m_n^{-1}$, independently of any structural player or season effects. An observation with 600 minutes of playing time is therefore inherently noisier than one with 2{,}500 minutes, even if both belong to players with identical latent profiles.

Left unaddressed, this heteroscedasticity causes the model to attribute the elevated variance of low-minute observations to heavy residual tails (via a smaller $\nu$) instead of to exposure, which is a misspecification. The production model removes it by scaling the residual variance at the observation level, as follows.

\begin{definition}[Minutes-weighted residual scale]
\label{def:phi}
Let $m_{\mathrm{ref}} > 0$ denote a reference exposure level, taken as the sample median of the minutes-played values $\{m_n\}_{n=1}^N$. Define the observation-level residual standard deviation for observation $n$ and feature $p$ as
\begin{equation}
\sigma_{e,n,p} \;=\; \sigma_{e,p} \cdot \left(\frac{m_{\mathrm{ref}}}{m_n}\right)^{\!\phi},
\label{eq:phi_scaling}
\end{equation}
where $\phi \geq 0$ is the scaling exponent. The Student-$t$ model with minutes-weighting replaces $\sigma_{e,p}$ in \eqref{eq:t_resid} with $\sigma_{e,n,p}$. This is the residual specification used by the production model throughout the remainder of this chapter and in Chapter~\ref{chap:results}: Definition~\ref{def:student_t} combined with \eqref{eq:phi_scaling}, not Definition~\ref{def:student_t} in isolation.
\end{definition}

\begin{corollary}[Variance under the minutes-weighted scale]
\label{cor:t_variance_mv}
Replacing $\sigma_{e,p}$ by $\sigma_{e,n,p}$ in Proposition~\ref{prop:t_variance} gives
\begin{equation}
\mathrm{Var}(\varepsilon_{np}) \;=\; \sigma_{e,n,p}^2\,\frac{\nu}{\nu-2} \;=\; \sigma_{e,p}^2 \left(\frac{m_{\mathrm{ref}}}{m_n}\right)^{\!2\phi}\frac{\nu}{\nu-2},
\label{eq:t_variance_mv}
\end{equation}
which depends on the observation $n$ through $m_n$ and reduces to Proposition~\ref{prop:t_variance}'s expression when $m_n = m_{\mathrm{ref}}$.
\end{corollary}

Three canonical values of $\phi$ are immediately interpretable. Setting $\phi = 0$ recovers the constant-scale case of \eqref{eq:t_resid}. Setting $\phi = 1/2$ corresponds to a square-root scaling analogous to the standard deviation of a Poisson count: variance grows linearly with event rate and standard deviation grows as the square root of inverse minutes. Setting $\phi = 1$ imposes full inverse-minutes scaling, so that the residual standard deviation is proportional to $m_n^{-1/2}$, as it would be for a simple average estimator under a constant-intensity Poisson process.

\begin{remark}[$\phi$ is fixed, not estimated]
\label{rem:phi}
The production model fixes $\phi = 1/2$, the Poisson-rate value described above, rather than treating it as a free parameter. Fixing $\phi$ introduces no additional parameter and requires no prior on $\phi$ in \eqref{eq:joint_prior}. The alternative --- assigning $\phi$ a prior and sampling it jointly with the remaining parameters --- is a well-posed model: identifiability of a free $\phi$ rests on the contrast between the residuals of high- and low-minute observations, since, given estimates of $\mathbf{a}_i$ and $\mathbf{b}_j$, the relationship between $|\varepsilon_{np}|$ and $m_n$ is informative about $\phi$, though $\phi$ must be identified jointly with $\sigma_{e,p}$ and $\nu$, all three of which govern the scale and shape of the marginal residual distribution. This alternative was fit as a diagnostic check, with $\phi$ assigned a $\mathrm{Normal}(0.5,\, 0.2)$ prior (Chapter~\ref{chap:results}): its posterior concentrates below $0.5$ but close to it, which is read there as support for fixing $\phi$ at the theoretically motivated value rather than as a reason to prefer the estimated one.
\end{remark}

\section{Model identification}
\label{sec:identification}

The factor model \eqref{eq:sigma_a} introduces a structural non-identifiability in the loading matrix $\boldsymbol{\Lambda}$ that must be resolved before Bayesian inference is defined.

\subsection{Rotational non-identifiability}

\begin{proposition}[Rotation equivalence]
\label{prop:rotation}
For any orthogonal matrix $\mathbf{Q} \in \mathbb{R}^{K \times K}$ satisfying $\mathbf{Q}\mathbf{Q}^\top = \mathbf{I}_K$, the loading matrices $\boldsymbol{\Lambda}$ and $\boldsymbol{\Lambda}\mathbf{Q}$ generate the same covariance:
\begin{equation}
(\boldsymbol{\Lambda}\mathbf{Q})(\boldsymbol{\Lambda}\mathbf{Q})^\top = \boldsymbol{\Lambda}\mathbf{Q}\mathbf{Q}^\top\boldsymbol{\Lambda}^\top = \boldsymbol{\Lambda}\boldsymbol{\Lambda}^\top,
\label{eq:rotation_equiv}
\end{equation}
and consequently $(\boldsymbol{\Lambda}\mathbf{Q})(\boldsymbol{\Lambda}\mathbf{Q})^\top + \boldsymbol{\Psi}_a = \boldsymbol{\Sigma}_a$.
\end{proposition}
\begin{proof}
Direct expansion using $\mathbf{Q}\mathbf{Q}^\top = \mathbf{I}_K$.
\end{proof}

Proposition~\ref{prop:rotation} implies that the map $(\boldsymbol{\Lambda}, \boldsymbol{\Psi}_a) \mapsto \boldsymbol{\Sigma}_a$ is many-to-one: any element of the orbit $\{\boldsymbol{\Lambda}\mathbf{Q} : \mathbf{Q} \in O(K)\}$ produces the same covariance, where $O(K)$ denotes the group of $K \times K$ orthogonal matrices. The group $O(K)$ has dimension $K(K-1)/2$, corresponding to the degrees of freedom of a $K$-dimensional rotation; it also includes reflections, so that there are $2^K$ discrete sign symmetries in addition to the continuous rotational freedom. Without constraints, the posterior distribution over $\boldsymbol{\Lambda}$ is multimodal and the sampler cannot mix across symmetry-equivalent modes.

\subsection{Lower-triangular positive-diagonal parametrisation}
\label{sec:ltpd}

A classical and computationally convenient solution to the non-identifiability described above is to impose a constraint on $\boldsymbol{\Lambda}$ that selects exactly one representative from each equivalence class under $O(K)$.

\begin{definition}[LT-PD constraint]
\label{def:ltpd}
A matrix $\boldsymbol{\Lambda} \in \mathbb{R}^{P \times K}$ satisfies the \emph{lower-triangular positive-diagonal} (LT-PD) constraint if:
\begin{enumerate}
\item[(i)] $\Lambda_{pk} = 0$ for all $p < k$ \quad (lower-triangular structure), and
\item[(ii)] $\Lambda_{kk} > 0$ for all $k = 1, \ldots, K$ \quad (positive diagonal).
\end{enumerate}
\end{definition}

\begin{proposition}[Identification via LT-PD]
\label{prop:ltpd_ident}
Suppose $P \geq K$ and let $\mathcal{L}$ denote the set of $P \times K$ matrices satisfying the LT-PD constraint. Then the map
\[
\mathcal{L} \times \{\boldsymbol{\Psi}_a : \psi_{a,p} > 0\;\forall\, p\} \;\longrightarrow\; \mathcal{S}^+_P, \qquad (\boldsymbol{\Lambda},\boldsymbol{\Psi}_a) \mapsto \boldsymbol{\Lambda}\boldsymbol{\Lambda}^\top + \boldsymbol{\Psi}_a,
\]
is injective, where $\mathcal{S}^+_P$ denotes the cone of $P \times P$ positive definite matrices. The factor covariance model is therefore identified under the LT-PD constraint.
\end{proposition}

The result in Proposition~\ref{prop:ltpd_ident} is established in \cite{anderson1956statistical}; the lower-triangular structure eliminates the continuous rotational freedom by zeroing out the $K(K-1)/2$ upper-triangular entries, while the positive-diagonal constraint eliminates the $2^K$ discrete sign symmetries. An accessible treatment in the Bayesian factor analysis context is given in \cite{geweke1996measuring}.

\begin{remark}[The LT-PD constraint is not an interpretive rotation]
\label{rem:ltpd_notinterp}
The LT-PD constraint resolves the mathematical non-identifiability \cite{anderson1956statistical,geweke1996measuring}, but it does not produce loadings that are interpretable in a football sense. A simple-structure rotation such as varimax or oblimin minimises cross-loadings and maximises the separation of features across factors, yielding loading patterns that correspond to recognisable dimensions. The LT-PD constraint does neither: it fixes a coordinate system by zeroing out the upper triangle of $\boldsymbol{\Lambda}$, a choice determined by feature ordering rather than by football logic. The resulting columns of $\boldsymbol{\Lambda}$ do not correspond to canonical latent directions and should not be reported or interpreted as such. Interpretable factor directions are recovered through the post-processing procedure described in Section~\ref{sec:spectral}.
\end{remark}

\subsection{Identified summaries of the factor structure}
\label{sec:spectral}

We report the spectral decomposition of the common variance matrix $\mathbf{C} = \boldsymbol{\Lambda}\boldsymbol{\Lambda}^\top$ rather than the raw loadings $\boldsymbol{\Lambda}$ themselves, because $\mathbf{C}$ is invariant to the LT-PD identification constraint (Proposition~\ref{prop:spectral_inv} below): it is uniquely determined by the model regardless of which constraint was used to fit $\boldsymbol{\Lambda}$. What follows are therefore identified summaries of the factor structure, not an interpretive rotation.

\begin{definition}[Spectral post-processing]
\label{def:pca_postproc}
For each posterior draw $s = 1, \ldots, S_{\mathrm{mc}}$, let $\boldsymbol{\Lambda}^{(s)}$ denote the sampled loading matrix. Compute the common variance matrix
\begin{equation}
\mathbf{C}^{(s)} = \boldsymbol{\Lambda}^{(s)}\!\left(\boldsymbol{\Lambda}^{(s)}\right)^\top \in \mathcal{S}^+_P,
\label{eq:common_variance}
\end{equation}
and perform its spectral decomposition
\begin{equation}
\mathbf{C}^{(s)} = \mathbf{V}^{(s)}\mathbf{D}^{(s)}\!\left(\mathbf{V}^{(s)}\right)^\top,
\label{eq:spectral_eq}
\end{equation}
where $\mathbf{D}^{(s)} = \mathrm{diag}\!\left(d_1^{(s)} \geq \cdots \geq d_K^{(s)} \geq 0\right)$ and $\mathbf{V}^{(s)} \in \mathbb{R}^{P \times K}$ has orthonormal columns. The $k$-th column $\mathbf{v}_k^{(s)}$ of $\mathbf{V}^{(s)}$ is the $k$-th principal eigenvector of $\mathbf{C}^{(s)}$; the corresponding eigenvalue $d_k^{(s)}$ is the amount of player-effect common variance explained by the $k$-th canonical factor direction. The rotated factor score for player $i$ in draw $s$ is
\begin{equation}
\mathbf{f}_i^{(s)} \;=\; \left(\mathbf{V}^{(s)}\right)^\top \boldsymbol{\eta}_i^{(s)} \;\in\; \mathbb{R}^K.
\label{eq:rotated_scores}
\end{equation}
\end{definition}

\begin{proposition}[Invariance of the spectral decomposition]
\label{prop:spectral_inv}
The eigenvalues $\{d_k^{(s)}\}_{k=1}^K$ and the common variance matrix $\mathbf{C}^{(s)}$ are invariant to any orthogonal transformation of the loading matrix. Replacing $\boldsymbol{\Lambda}^{(s)}$ by $\boldsymbol{\Lambda}^{(s)}\mathbf{Q}$ for any $\mathbf{Q} \in O(K)$ yields the same $\mathbf{C}^{(s)}$ and hence the same spectral decomposition \eqref{eq:spectral_eq}.
\end{proposition}
\begin{proof}
By Proposition~\ref{prop:rotation}, $(\boldsymbol{\Lambda}^{(s)}\mathbf{Q})(\boldsymbol{\Lambda}^{(s)}\mathbf{Q})^\top = \boldsymbol{\Lambda}^{(s)}\boldsymbol{\Lambda}^{(s)\top} = \mathbf{C}^{(s)}$.
\end{proof}

\begin{remark}[Sign convention]
\label{rem:sign}
Eigenvectors are defined up to sign: if $\mathbf{v}$ is an eigenvector of $\mathbf{C}^{(s)}$, so is $-\mathbf{v}$. A sign convention is applied draw-by-draw before computing posterior summaries: for each canonical factor $k$, the sign of $\mathbf{v}_k^{(s)}$ is chosen so that the feature with the largest absolute loading on that eigenvector has a positive loading. This convention ensures that the posterior mean of $\mathbf{v}_k^{(s)}$ reflects a coherent direction rather than cancelling between draws of opposite sign.
\end{remark}

\section{Prior specification}
\label{sec:priors}

All model parameters are assigned priors that are weakly informative relative to the scale of the Z-standardised outcome $\mathbf{y}_n$. Because the features have unit variance by construction, a loading of magnitude $|\Lambda_{pk}| = 1$ represents a one-standard-deviation shift in feature $p$ per unit change in the $k$-th factor score, which is already a large effect for a Z-scored performance variable. Priors concentrated near zero with tails extending to values of order one are therefore practically and scientifically reasonable.

\begin{definition}[Priors on the loading matrix]
\label{def:prior_lambda}
The diagonal entries of $\boldsymbol{\Lambda}$ are assigned
\begin{equation}
\Lambda_{kk} \;\sim\; \mathrm{Half\text{-}Normal}(0,\, 1), \qquad k = 1, \ldots, K,
\label{eq:prior_lambda_diag}
\end{equation}
consistent with the LT-PD positivity constraint. The free lower-triangular entries are assigned
\begin{equation}
\Lambda_{pk} \;\sim\; \mathcal{N}(0,\, 1), \qquad 1 \leq k < p \leq P.
\label{eq:prior_lambda_off}
\end{equation}
\end{definition}

\begin{definition}[Priors on variance parameters]
\label{def:prior_variances}
The uniqueness standard deviations, the season standard deviations, and the residual standard deviations are each assigned independent Half-Normal priors:
\begin{align}
\psi_{a,p} &\;\sim\; \mathrm{Half\text{-}Normal}(0,\, 1), \notag \\
\sigma_{b,p} &\;\sim\; \mathrm{Half\text{-}Normal}(0,\, 1), \notag \\
\sigma_{e,p} &\;\sim\; \mathrm{Half\text{-}Normal}(0,\, 1), \qquad p = 1, \ldots, P.
\label{eq:prior_variances}
\end{align}
\end{definition}

The Half-Normal$(0,1)$ prior for scale parameters concentrates most mass near zero while permitting values of order one with positive probability; it is recommended for variance components in hierarchical models as a weakly informative alternative to flat or improper priors, which can produce improper posteriors in hierarchical settings \cite{gelman2006prior}.

\begin{definition}[Prior on the degrees-of-freedom parameter]
\label{def:prior_nu}
The Student-$t$ degrees-of-freedom parameter is assigned
\begin{equation}
\nu \;\sim\; \mathrm{Gamma}(2,\, 0.1),
\label{eq:prior_nu}
\end{equation}
supported on $\nu > 2$ to ensure finite residual variance. The $\mathrm{Gamma}(2, 0.1)$ distribution has mean $20$ and standard deviation approximately $14$. It places substantial mass below $\nu = 10$, heavy tails, while also allowing the posterior to concentrate near larger values if the data support behaviour close to Gaussian. The prior does not favour either the Gaussian limit $\nu \to \infty$ nor the extreme heavy-tail regime $\nu \to 2^+$\cite{juarez2010model}.
\end{definition}

Since the production model fixes $\phi=1/2$ (Remark~\ref{rem:phi}) rather than estimating it, $\phi$ is not a random variable and is not assigned a prior here; it enters \eqref{eq:phi_scaling} as a known constant. (The diagnostic fit mentioned in Remark~\ref{rem:phi}, which does estimate $\phi$, assigns it $\phi \sim \mathrm{Normal}(0.5,\, 0.2)$ --- reported there for completeness, not as part of the production prior stack below.)

The joint prior is the product of all terms in Definitions~\ref{def:prior_lambda}--\ref{def:prior_nu}, since all parameters are treated as independent then the joint prior density is:
\begin{align}
p\!\left(\boldsymbol{\Lambda}, \boldsymbol{\psi}_a, \boldsymbol{\sigma}_b, \boldsymbol{\sigma}_e, \nu\right)
\;=\;
&\prod_{k=1}^{K} p_{\mathrm{HN}}(\Lambda_{kk})
\prod_{\substack{p=1\\p>k}}^{P} \mathcal{N}(\Lambda_{pk};\, 0, 1) \notag \\
&\;\times\;
\prod_{p=1}^{P} p_{\mathrm{HN}}(\psi_{a,p})\,p_{\mathrm{HN}}(\sigma_{b,p})\,p_{\mathrm{HN}}(\sigma_{e,p}) \notag \\
&\;\times\;
p_{\mathrm{Ga}}(\nu),
\label{eq:joint_prior}
\end{align}
where $p_{\mathrm{HN}}$ denotes the Half-Normal$(0,1)$ density and $p_{\mathrm{Ga}}$ denotes the Gamma$(2,0.1)$ density; $\phi$ does not appear, being fixed rather than estimated.

\section{Rank selection}
\label{sec:rank}

The factor rank $K$ governs the dimension of the latent space used to model the player covariance $\boldsymbol{\Sigma}_a$. Unlike the parameters $\boldsymbol{\Lambda}$, $\boldsymbol{\Psi}_a$, and $\nu$, $K$ could in principle be given a prior and estimated jointly with the rest of the model; here it is instead fixed by a qualitative, post-hoc argument, since $K$ determines the parametric family itself and a within-model treatment was not pursued.

For any fixed value of $K \in \{0, 1, 2, \ldots\}$, the model specification in Sections~\ref{sec:covariance}--\ref{sec:priors} is complete. The case $K = 0$ is a special case: it sets $\boldsymbol{\Sigma}_a = \boldsymbol{\Psi}_a$ (diagonal player covariance), corresponding to the assumption that all player-level cross-feature correlation is zero. This \emph{diagonal model} serves as a natural baseline; the gain in predictive performance relative to it quantifies the value of modelling shared player-level covariance structure.

Whether shared player-level covariance is worth modelling at all is answered by the posterior-predictive check reported in Chapter~\ref{chap:results} (\S\ref{sec:results-rank}), which compares model-implied and observed cross-feature correlations under $K=0$ and $K\geq1$ directly, without appeal to cross-validation; the choice of rank among $K \geq 1$ is argued qualitatively, as follows.

\subsection{A qualitative argument for the rank}
\label{sec:rank_qualitative}

One caveat is worth flagging briefly before making that argument: the rank $K^*=3$ selected here is not expected to match the eleven-factor count found in Chapter~\ref{chap:background}'s exploratory analysis. The two describe different things.

The rank actually adopted for the production model is selected on qualitative, structural grounds instead: the adequacy of each fit's own residuals, and the extent to which the shared player-effect covariance $\boldsymbol{\Sigma}_a$ is still adding well-identified structure as $K$ increases. This mirrors standard practice in exploratory and confirmatory factor analysis, where the number of retained factors is routinely set by fit diagnostics and interpretability rather than by a single automated rule, and is consistent with the professional decision making process in my day-to-day job: that a qualitative argument, grounded in residuals and in the variance structure of $\boldsymbol{\Lambda}\boldsymbol{\Lambda}^\top$, is an appropriate and sufficient basis for closing the choice of $K$.

\begin{remark}[Anchoring to the posterior mean, not to per-draw estimates]
\label{rem:posterior_mean_anchor}
Both criteria below require the eigendecomposition of the common-variance matrix $\boldsymbol{\Lambda}\boldsymbol{\Lambda}^\top$ introduced in Section~\ref{sec:spectral}. Proposition~\ref{prop:spectral_inv} shows that, \emph{for a fixed posterior draw}, this eigendecomposition is invariant to the identification constraint used to fit $\boldsymbol{\Lambda}$. It is not, however, invariant \emph{across} posterior draws in the sense needed to build a credible interval directly: sampling variability in $\boldsymbol{\Lambda}^{(s)}$ can rotate the leading eigenvectors within the subspace they span from one draw to the next, so naively comparing $\mathbf{V}^{(s)}$ across draws risks comparing different rotations of the same subspace rather than a stable estimate of one direction. To avoid this, the analyses below fix a single reference orientation---the eigendecomposition of the \emph{posterior-mean} common-variance matrix $\bar{\boldsymbol{\Lambda}}\bar{\boldsymbol{\Lambda}}^\top$, where $\bar{\boldsymbol{\Lambda}} = \tfrac{1}{S_{\mathrm{mc}}}\sum_s \boldsymbol{\Lambda}^{(s)}$---and project each draw's loadings onto this fixed reference before summarising uncertainty. The same construction gives comparable per-player credible intervals: replacing $\mathbf{V}^{(s)}$ by the fixed $\bar{\mathbf{V}}$ in \eqref{eq:rotated_scores} yields $\mathbf{f}_i^{(s)} = \bar{\mathbf{V}}^\top\boldsymbol{\Lambda}^{(s)}\boldsymbol{\eta}_i^{(s)}$, used in Chapter~\ref{chap:results} to report posterior credible intervals on individual players' canonical scores.
\end{remark}

\begin{description}
\item[Criterion A: variance share of the PCs of $\bar{\boldsymbol{\Lambda}}\bar{\boldsymbol{\Lambda}}^\top$] For a candidate rank $K$, the eigenvalue $d_K$ of the \emph{last} (smallest) canonical factor, expressed as a share of $\sum_{k=1}^K d_k$, measures how much additional common-variance structure the $K$-th dimension contributes. A last-component share close to the trailing eigenvalues of the full uniqueness spectrum indicates that dimension is not resolving substantive shared structure.
\item[Criterion B: residual adequacy under the fit's own $t(\hat\nu)$ noise model] For each candidate fit, the standardised residuals $z_{np} = (y_{np} - \hat a_{i[n],p} - \hat b_{j[n],p})/(\hat\sigma_{e,n,p}\sqrt{\hat\nu/(\hat\nu-2)})$ are compared against the $99$th-percentile threshold $\tau_{99}$ implied by $t(\hat\nu)$, where $\hat\sigma_{e,n,p}$ is the fit's own residual scale at observation $n$ (Definition~\ref{def:phi} for a minutes-scaled fit; $\hat\sigma_{e,n,p}=\hat\sigma_{e,p}$, constant in $n$, for a constant-scale fit). If the empirical exceedance rate across all $N \times P$ cells is at or below the nominal $1\%$, the residual distribution is adequate under the fitted noise model; this is a floor check on model misspecification, not a criterion that discriminates between two ranks that both pass it.
\end{description}

\subsection{Selected rank}
\label{sec:rank_selection}

Applying Criteria A and B across $K=2, 3, 4$ gives the following picture. The last canonical factor's common-variance share declines steadily from $K=2$ to $K=3$ to $K=4$---a continued but decelerating decline, not a sharp drop to a noise floor at any point (Criterion A). Mean residual exceedance under $\tau_{99}$ stays below the $1\%$ floor and essentially flat across all three ranks, so residual adequacy holds at every rank examined and does not discriminate among them (Criterion B).

Taken together, this comparison does not produce a clean reversal at $K=4$ the way it does at $K=3$ relative to $K=2$: Criterion A mildly favours stopping at $K=3$ (the fourth factor's contribution, while still non-trivial, is smaller than the third's), while Criterion B is uninformative at this rank range and offers no separate support either way. Choosing $K=3$ over $K=4$ rests on the preference for the simpler model. This is a much weaker result than the case for $K\geq 1$ over $K=0$, or even the milder case for $K=3$ over $K=2$; both of those were clear-cut, this one is not. The honest summary is that the data do not sharply distinguish $K=3$ from $K=4$. $K=3$ is nonetheless retained as the production rank $K^*=3$, used throughout the remainder of this chapter and in Chapter~\ref{chap:results}: it is the more parsimonious of two ranks that the evidence does not clearly separate, and closing the model at a defensible, adequately-supported rank was judged preferable to an open-ended search for further, likely similarly ambiguous, comparisons at $K \geq 5$.


\section{Variance decomposition and intraclass correlation}
\label{sec:icc}

A key inferential output of the model is a decomposition of the total variance of each feature into contributions from the player, season, and residual components. Under the additive model \eqref{eq:additive} and the mutual independence of $\mathbf{a}_i$, $\mathbf{b}_j$, and $\boldsymbol{\varepsilon}_n$, the marginal variance of feature $p$ in observation $n$ is
\begin{equation}
\mathrm{Var}(y_{np}) \;=\; [\boldsymbol{\Sigma}_a]_{pp} + \sigma_{b,p}^2 + \mathrm{Var}(\varepsilon_{np}),
\label{eq:total_variance}
\end{equation}
where $[\boldsymbol{\Sigma}_a]_{pp}$ is given by \eqref{eq:sigma_diag}, and, for the production model's minutes-weighted Student-$t$ residual, $\mathrm{Var}(\varepsilon_{np}) = \sigma_{e,n,p}^2 \cdot \nu/(\nu-2)$ by Corollary~\ref{cor:t_variance_mv} --- a quantity that depends on $n$ through the observation's playing time $m_n$, not on $p$ alone.

Because $\mathrm{Var}(\varepsilon_{np})$ varies across observations of the same feature, so does $\mathrm{Var}(y_{np})$: there is no longer a single total variance, or a single variance decomposition, per feature, only one per feature at each exposure level. The definition below resolves this by fixing $m_n = m_{\mathrm{ref}}$, the reference exposure of Definition~\ref{def:phi}, at which the scaling factor $(m_{\mathrm{ref}}/m_n)^\phi$ equals one and $\sigma_{e,n,p}$ reduces to $\sigma_{e,p}$. This recovers a single number per feature, interpretable as the variance decomposition for a player--season observation at the reference playing time.

\begin{definition}[Intraclass correlation coefficients]
\label{def:icc}
The player, season, and residual intraclass correlation coefficients for feature $p$, evaluated at the reference exposure $m_n = m_{\mathrm{ref}}$, are, respectively,
\begin{align}
\mathrm{ICC}_{a,p} &= \frac{[\boldsymbol{\Sigma}_a]_{pp}}{[\boldsymbol{\Sigma}_a]_{pp} + \sigma_{b,p}^2 + \mathrm{Var}(\varepsilon_{np})},
\label{eq:icc_a} \\[4pt]
\mathrm{ICC}_{b,p} &= \frac{\sigma_{b,p}^2}{[\boldsymbol{\Sigma}_a]_{pp} + \sigma_{b,p}^2 + \mathrm{Var}(\varepsilon_{np})},
\label{eq:icc_b} \\[4pt]
\mathrm{ICC}_{e,p} &= \frac{\mathrm{Var}(\varepsilon_{np})}{[\boldsymbol{\Sigma}_a]_{pp} + \sigma_{b,p}^2 + \mathrm{Var}(\varepsilon_{np})},
\label{eq:icc_e}
\end{align}
satisfying $\mathrm{ICC}_{a,p} + \mathrm{ICC}_{b,p} + \mathrm{ICC}_{e,p} = 1$ for all $p$.
\end{definition}

The quantity $\mathrm{ICC}_{a,p}$ measures the proportion of total variance in feature $p$ that is attributable to stable differences among players. It is directly interpretable as the reliability of the feature as a measure of stable player style: features with $\mathrm{ICC}_{a,p}$ close to one are dominated by stable player characteristics, while features with small $\mathrm{ICC}_{a,p}$ are dominated by season-level context or observation-level noise. This framing of variance decomposition in terms of intraclass correlation follows \cite{shrout1979intraclass}.

\begin{remark}[Bayesian ICC]
\label{rem:bayesian_icc}
Because all variance parameters are sampled from the posterior, the quantities defined in Definition~\ref{def:icc} are themselves posterior quantities. For each posterior draw, a value of $\mathrm{ICC}_{a,p}$ can be computed from the sampled values of $\boldsymbol{\Lambda}$, $\boldsymbol{\Psi}_a$, $\sigma_{b,p}$, $\sigma_{e,p}$, and $\nu$, yielding a full posterior distribution over each coefficient. Posterior medians and $95\%$ credible intervals are reported in Chapter~\ref{chap:results}. The ICC framing is preferred over the raw variance decomposition because it is scale-free and comparable across features with different marginal variances.
\end{remark}


